\documentclass[11pt]{article}
\usepackage{amsmath,amssymb,amsthm}
\usepackage{geometry}
\geometry{margin=1in}
\usepackage{hyperref}
\hypersetup{colorlinks=true,linkcolor=blue,citecolor=blue,urlcolor=blue}
\usepackage[numbers]{natbib}
\usepackage{booktabs}
\usepackage{parskip}

\title{The moment decomposition for time-varying \\
       Johnson--Mehl--Avrami--Kolmogorov kinetics}
\author{Documentation for the JMAK per-voxel machinery in THAMES \\
        \texttt{namespace jmak} in \texttt{JMAKGrowth.h/.cc}}
\date{28 July 2026}

\begin{document}
\maketitle

\section{Purpose}

This note explains how the four moment accumulators
$M_0, M_1, M_2, M_3$ used in the THAMES JMAK code
(\texttt{jmak::GlobalMoments}) reduce the doubly-nested time integral
of extended volume to a per-cycle constant-time update, when both
the nucleation rate $J(t)$ and the linear growth velocity $G(t)$
depend on time.
The classical constant-$J$, constant-$G$ result is assumed familiar;
we build on it directly and derive the moment identity used in the
implementation.

The technique complements two treatments the reader may already know:
Scherer and Bellmann\cite{Scherer2018} handle time-varying driving
force via an \emph{effective time} variable
$\xi(t)$; Bullard, Scherer, and Thomas\cite{Bullard2015} handle it in
HydratiCA by discretizing rate laws numerically at each time step.
The moment decomposition is a third route: it is algebraic (no change
of variable, no numerical quadrature over history) and matches the
structure of the JMAK integrand exactly for integer $n=4$
three-dimensional continuous nucleation, which is the case implemented
in the current THAMES code.

\section{Classical JMAK, briefly}

For constant $J$ (events per unit volume per unit time) and constant
$G$ (linear growth velocity of the nucleus--matrix interface), the
Johnson--Mehl--Avrami--Kolmogorov result for the transformed volume
fraction in a bulk sample is
\begin{equation}
X(t) \;=\; 1 - \exp\!\left(-K t^{n}\right)\,,
\label{eq:classical-jmak}
\end{equation}
where $K = (\pi / 3)\, J\, G^{3}$ and $n = 4$ for three-dimensional
isotropic growth with continuous (site-nonsaturating) nucleation;
$n = 3$ if all nuclei are present at $t=0$ (site-saturated).
Equation~\eqref{eq:classical-jmak} follows from the observation that
the naive
``extended'' volume fraction $Y_e(t)$, ignoring impingement between
growing particles, is
\begin{equation}
Y_e(t) \;=\; \frac{4\pi}{3}\int_{0}^{t} J\,\bigl[G\,(t-\tau)\bigr]^{3}\,\mathrm{d}\tau
       \;=\; \frac{\pi}{3}\, J\, G^{3}\, t^{4}\,,
\label{eq:Ye-constant-JG}
\end{equation}
and the transformation from extended to actual is
\begin{equation}
X(t) \;=\; 1 - \exp[-Y_e(t)]\,,
\label{eq:X-from-Ye}
\end{equation}
Kolmogorov's Poisson-impingement correction.
The integrand of \eqref{eq:Ye-constant-JG} is
\emph{(events per unit volume per unit time)} $\times$
\emph{(volume of a spherical nucleus that has grown for
$t - \tau$)}, integrated over the birth time $\tau$.

\section{Time-varying $J$ and $G$}

When the driving force changes with time, both $J$ and $G$ become
functions of $t$ (through the current saturation index or
supersaturation, which in cement systems is governed by the coupled
evolution of the dissolving reactants and the solid product).
The nucleus born at $\tau$ then has, at time $t > \tau$, a radius
\begin{equation}
r(\tau, t) \;=\; \int_{\tau}^{t} G(s)\, \mathrm{d}s\,,
\label{eq:r-tau-t}
\end{equation}
which is not simply $G\,(t-\tau)$ because $G$ varies over the growth
interval.
The extended volume in a fixed system volume becomes the doubly-nested
integral
\begin{equation}
Y_e(t) \;=\; \frac{4\pi}{3}\int_{0}^{t}\!\! J(\tau)\,
    \Bigl[\underbrace{\int_{\tau}^{t}\! G(s)\, \mathrm{d}s}_{r(\tau, t)}\Bigr]^{3}
    \mathrm{d}\tau\,.
\label{eq:Ye-varying}
\end{equation}
Equation~\eqref{eq:Ye-varying} is what a naive discretization would
carry: at each cycle $t_k$ we would need to evaluate the outer
integral by iterating over every past cycle, and for each we would
evaluate the inner integral over the interval $[\tau, t_k]$.
The cost per cycle is proportional to the number of past cycles, so
total work over $N$ cycles is $O(N^2)$.
For $N \sim 10^4$ cycles this is $10^8$ operations per phase per
cycle; expensive but not prohibitive.
The moment identity below reduces it to $O(1)$ per cycle per phase
without approximation, which matters when the code has multiple
phases and the run is long.

\section{Change of accumulator}

Introduce the cumulative growth arc
\begin{equation}
r_{\text{acc}}(t) \;=\; \int_{0}^{t}\! G(s)\, \mathrm{d}s\,.
\label{eq:racc-def}
\end{equation}
This is a length, not a rate; it is the radius that a hypothetical
nucleus born at $t = 0$ would have reached by time $t$
if $G$ had continued at its actual (possibly time-varying) values.
For a nucleus actually born at $\tau > 0$, its radius at time $t$
follows immediately from~\eqref{eq:r-tau-t}:
\begin{equation}
r(\tau, t) \;=\; r_{\text{acc}}(t) \;-\; r_{\text{acc}}(\tau)\,.
\label{eq:r-tau-t-diff}
\end{equation}
The \emph{difference} of the same accumulator, evaluated at the
current time and at the birth time.
This is the key.

Substitute~\eqref{eq:r-tau-t-diff} into~\eqref{eq:Ye-varying}:
\begin{equation}
Y_e(t) \;=\; \frac{4\pi}{3}\int_{0}^{t}\! J(\tau)\,
    \bigl[r_{\text{acc}}(t) - r_{\text{acc}}(\tau)\bigr]^{3}\, \mathrm{d}\tau\,.
\label{eq:Ye-with-racc}
\end{equation}
The cube in the integrand can now be expanded by the binomial theorem.
The term $r_{\text{acc}}(t)$ is independent of the integration variable
$\tau$, so it can be pulled out of the integral, exposing the
$\tau$-dependence entirely in powers of $r_{\text{acc}}(\tau)$.

\section{Moment decomposition}

Expanding the cube,
\begin{align}
\bigl[r_{\text{acc}}(t) - r_{\text{acc}}(\tau)\bigr]^{3} &=
     r_{\text{acc}}(t)^{3} - 3\, r_{\text{acc}}(t)^{2}\, r_{\text{acc}}(\tau) \notag \\
&\quad + 3\, r_{\text{acc}}(t)\, r_{\text{acc}}(\tau)^{2} - r_{\text{acc}}(\tau)^{3}\,.
\end{align}
Substituting into~\eqref{eq:Ye-with-racc} and separating
$t$-dependent factors from $\tau$-dependent factors,
\begin{align}
Y_e(t) &= \frac{4\pi}{3}\left[
    r_{\text{acc}}(t)^{3}\, \int_{0}^{t}\! J(\tau)\, \mathrm{d}\tau
  - 3\, r_{\text{acc}}(t)^{2}\, \int_{0}^{t}\! J(\tau)\, r_{\text{acc}}(\tau)\, \mathrm{d}\tau
  \right. \notag \\
&\hphantom{= \frac{4\pi}{3}\left[\right.}\left.
  + 3\, r_{\text{acc}}(t)\,     \int_{0}^{t}\! J(\tau)\, r_{\text{acc}}(\tau)^{2}\, \mathrm{d}\tau
  -                              \int_{0}^{t}\! J(\tau)\, r_{\text{acc}}(\tau)^{3}\, \mathrm{d}\tau
  \right]\,.
\label{eq:Ye-decomposed}
\end{align}
Define the four \emph{moments} of the accumulated growth radius:
\begin{equation}
M_{k}(t) \;=\; \int_{0}^{t}\! J(\tau)\, r_{\text{acc}}(\tau)^{k}\, \mathrm{d}\tau\,,
\qquad k \in \{0, 1, 2, 3\}\,.
\label{eq:moments-def}
\end{equation}
$M_k(t)$ is the number-density-weighted $k$-th moment of the growth
radius at time $t$: how many nuclei have been produced up to $t$,
weighted by the $k$-th power of the accumulated growth radius they
have seen since $t=0$.

With this definition,~\eqref{eq:Ye-decomposed} becomes
\begin{equation}
\boxed{\;
Y_e(t) \;=\; \frac{4\pi}{3}\left[
     r_{\text{acc}}(t)^{3}\, M_{0}(t)
  \;-\; 3\, r_{\text{acc}}(t)^{2}\, M_{1}(t)
  \;+\; 3\, r_{\text{acc}}(t)\, M_{2}(t)
  \;-\; M_{3}(t)
\right]\,.
\;}
\label{eq:Ye-moments}
\end{equation}
This is the identity implemented in
\texttt{jmak::extendedVolumePerVoxel} (with \texttt{alpha} for the
morphology coefficient in place of $4\pi/3$).
Each moment $M_k(t)$ satisfies the trivial update rule
\begin{equation}
M_k(t + \Delta t) \;=\; M_k(t) + J(t)\, r_{\text{acc}}(t)^{k}\, \Delta t\,,
\label{eq:moment-update}
\end{equation}
at first order in $\Delta t$, which is one addition and $k$
multiplications per cycle.
Similarly $r_{\text{acc}}(t + \Delta t) = r_{\text{acc}}(t) + G(t)\, \Delta t$.
Once a phase's $r_{\text{acc}}$ and $M_0, \dots, M_3$ are updated for the
current cycle, the extended volume is one arithmetic evaluation of
\eqref{eq:Ye-moments}: five scalars in, one scalar out, no history
lookup.

For a single sample of volume $V$, the transformed fraction is the
usual $X(t) = 1 - \exp[-Y_e(t)/V]$; in the THAMES per-voxel setting,
$V$ is the volume of a single lattice voxel and each seeded voxel
carries its own $M_k$ snapshot taken at seed time (its ``generation''
in the code) so that a voxel's own $Y_e$ starts from zero when it
seeds, not from $t = 0$.
The generation snapshot is subtracted:
$Y_e^{\text{gen}}(t) = Y_e(t; \text{using } M_k(t) - M_k(t_{\text{seed}}))$.

\section{Per-cycle numerical update}

Only five scalars are stored per phase, plus a small snapshot per
generation. The per-cycle work is:
\begin{itemize}
\item Evaluate $J(t)$ from the CNT rate law at current $S(t)$
      (\texttt{cnt::nucleationRate}).
\item Evaluate $G(t)$ from the phase's surface-normal rate law times
      the molar volume: $G = r_{\text{surface}}(S)\, v_{\text{molar}}$
      (\texttt{getGrowthVelocity} in each concrete
      \texttt{KineticModel} subclass).
\item Advance $r_{\text{acc}}$ by $G\, \Delta t$.
\item Advance each $M_k$ by $J\, r_{\text{acc}}^{k}\, \Delta t$.
\item For each generation, evaluate~\eqref{eq:Ye-moments} using the
      difference between the current $M_k$ and the generation's
      snapshot; convert to $X_{\text{gen}}$ via $-\text{expm1}(-Y_e/V)$.
\end{itemize}
None of these steps looks up the past values of $J$ or $G$.
The full growth history is encoded in $r_{\text{acc}}$ and the four
moments.

\section{Connection with Scherer's effective-time formulation}

Scherer and Bellmann\cite{Scherer2018} handle a time-varying driving
force by defining an effective time
\begin{equation}
\xi(t) \;=\; \int_{0}^{t} \bigl(\beta_{\text{Tot}}(t')^{1/m} - 1\bigr)^{n}\, \mathrm{d}t'\,,
\label{eq:scherer-xi}
\end{equation}
where $\beta_{\text{Tot}}$ is the total-saturation ratio.
In their treatment, the linear dimension of each particle is
$x_i = G_{0i}\, \xi$: the entire time evolution is absorbed into a
single monotone variable~$\xi$, and the classical JMAK expressions
for constant growth rate apply in the $\xi$ coordinate.
The two devices — Scherer's effective time and this note's moment
decomposition — are addressing the same computational problem
(algebraic simplification of a double integral with a time-varying
kernel) via different mathematical routes.
Scherer's change of variable is exact when the shape factor $g$
(their Eq.~8) and the ratio $G_{0i}/G_{0j}$ are constants, i.e., when
the driving-force time dependence is separable from the direction
dependence.
The moment decomposition presented here is exact when the growth
kernel $r(\tau, t) = r_{\text{acc}}(t) - r_{\text{acc}}(\tau)$
separates into $t$- and $\tau$-dependent pieces after binomial
expansion, which is possible for integer positive powers of the
kernel.
The two coincide in the special case where the whole time-dependence
of the rates is captured by a single $\xi$.

\section{Extensions and limitations}

\begin{itemize}
\item \textbf{Non-integer exponent $n$.} The moment identity
      \eqref{eq:Ye-moments} requires $[r_{\text{acc}}(t) -
      r_{\text{acc}}(\tau)]^{m}$ to be expandable as a finite sum of
      products of $t$-only and $\tau$-only factors. This is possible
      for integer $m$ (via the binomial theorem). For non-integer
      $m$ — physically defensible for non-spherical or fractal
      morphologies where $n \in [2.5, 4]$ — the identity does not
      close, and per-cohort numerical integration of the growth
      kernel is required. The THAMES code currently implements only
      $m = 3$ (equivalently $n = 4$); non-integer $n$ is planned as
      an extension.

\item \textbf{Site-saturated nucleation.} If all nuclei are placed
      instantaneously at $t = t_c$ and no further nucleation occurs,
      $J(t) = N_0\, \delta(t - t_c)$ and
      $M_k(t) = N_0\, r_{\text{acc}}(t_c)^{k}$ for $t \geq t_c$.
      The formulation still applies, and the classical result
      $Y_e = (4\pi/3)\, N_0\, r(t_c, t)^{3}$ is recovered.

\item \textbf{Non-spherical shapes and heterogeneous nucleation.} The
      morphology coefficient $\alpha$ replaces $4\pi/3$ for
      non-spherical growth (e.g., $\alpha \approx \pi g / 2$ for
      blade-like C-S-H per Scherer's shape factor). For heterogeneous
      nucleation on a scarce substrate the effective $J$ must be
      further weighted by the fraction of available substrate; the
      moment update still applies to the effective $J$, but the
      substrate-availability multiplier is a separate factor whose
      time dependence must be tracked. In THAMES, the CNT rate law
      \texttt{cnt::nucleationRate} already carries a $\theta$
      parameter for that (Turnbull--Vonnegut $f(\theta)$).

\item \textbf{Impingement between voxels.} JMAK's Poisson-impingement
      correction $X = 1 - \exp(-Y_e/V)$ operates \emph{within a
      voxel}. Impingement between neighboring lattice voxels is
      handled at the placement stage
      (\texttt{Lattice::nucleatePhaseRnd}): once a voxel is fully
      transformed ($X_{\text{gen}} \approx 1$), its Portlandite
      contribution is committed as a whole voxel of the phase and no
      further JMAK evolution applies to it. Growth into neighboring
      electrolyte is handled by the standard THAMES rate-law/
      lattice-face mechanics.
\end{itemize}

\begin{thebibliography}{9}
\bibitem{Scherer2018}
G.\,W. Scherer, F. Bellmann,
\emph{Kinetic analysis of C-S-H growth on calcite},
Cement and Concrete Research \textbf{103}, 226--235 (2018).

\bibitem{Bullard2015}
J.\,W. Bullard, G.\,W. Scherer, J.\,J. Thomas,
\emph{Time dependent driving forces and the kinetics of tricalcium
silicate hydration},
Cement and Concrete Research \textbf{74}, 26--34 (2015).

\bibitem{Christian2002}
J.\,W. Christian,
\emph{The Theory of Transformations in Metals and Alloys},
Pergamon (2002).

\bibitem{AvramovSestak}
I. Avramov, J. \v{S}est\'{a}k,
Generalized kinetics of non-isothermal transformation phenomena,
\emph{Journal of Thermal Analysis and Calorimetry} \textbf{79},
595--602 (2005).
\end{thebibliography}

\end{document}
